% =============================================================================
% Section 4: Domain Gap Analysis
% =============================================================================
\section{Domain Gap Analysis}
\label{sec:domain_gap}

BrainSegFounder was pretrained and fine-tuned on glioma data. Before adapting
it to the meningioma domain, we quantify the distributional shift in the
frozen encoder's feature space to establish the empirical necessity of domain
adaptation.

% ---------------------------------------------------------------------------
% 4.1 Feature Extraction
% ---------------------------------------------------------------------------
\subsection{Feature Extraction Protocol}
\label{sec:dgap:extraction}

Features are extracted from the frozen BrainSegFounder encoder using the
\texttt{FeatureExtractor} pipeline:
\begin{equation}
  \mathbf{h} = \operatorname{GAP}\bigl(\texttt{encoder10}\bigl(
    \text{SwinViT}(\mathbf{x})\bigr)\bigr) \in \R^{768},
  \label{eq:feature_extraction}
\end{equation}
where $\mathbf{x} \in \R^{4 \times 192^3}$ is the input volume and
$\operatorname{GAP}$ denotes Global Average Pooling over the $4^3$ (or $6^3$
for $192^3$ input) spatial dimensions. We extract features from 200~random
BraTS-GLI subjects and 200~BraTS-MEN subjects from the test and SDP-train
splits.

% ---------------------------------------------------------------------------
% 4.2 Metrics
% ---------------------------------------------------------------------------
\subsection{Domain Gap Metrics}
\label{sec:dgap:metrics}

\subsubsection{Maximum Mean Discrepancy}

The Maximum Mean Discrepancy (MMD) \cite{gretton2012kernel} measures the
distance between two distributions in a reproducing kernel Hilbert space.
Given samples $\{\mathbf{h}_i^{\text{GLI}}\}_{i=1}^{n}$ and
$\{\mathbf{h}_j^{\text{MEN}}\}_{j=1}^{m}$, the unbiased estimator is:
\begin{equation}
  \widehat{\text{MMD}}^2 = \frac{1}{n(n-1)}\sum_{i \neq i'} k(\mathbf{h}_i,
    \mathbf{h}_{i'}) + \frac{1}{m(m-1)}\sum_{j \neq j'} k(\mathbf{h}_j,
    \mathbf{h}_{j'}) - \frac{2}{nm}\sum_{i,j} k(\mathbf{h}_i, \mathbf{h}_j),
  \label{eq:mmd}
\end{equation}
where $k$ is a Gaussian kernel with bandwidth $\sigma$ set to the median
pairwise distance (the median heuristic). Statistical significance is
assessed via a permutation test with $n_{\text{perm}} = 1{,}000$
permutations.

\subsubsection{Centred Kernel Alignment}

CKA \cite{kornblith2019similarity} measures representational similarity
between two feature matrices $\mathbf{H}_1 \in \R^{n \times d_1}$ and
$\mathbf{H}_2 \in \R^{n \times d_2}$ via:
\begin{equation}
  \text{CKA}(\mathbf{H}_1, \mathbf{H}_2) = \frac{
    \|\bar{\mathbf{H}}_1^\top \bar{\mathbf{H}}_2\|_F^2
  }{
    \|\bar{\mathbf{H}}_1^\top \bar{\mathbf{H}}_1\|_F \cdot
    \|\bar{\mathbf{H}}_2^\top \bar{\mathbf{H}}_2\|_F
  },
  \label{eq:cka}
\end{equation}
where $\bar{\mathbf{H}}$ denotes column-centred matrices. CKA $\in [0, 1]$,
with 1 indicating identical representations (up to linear transformation).

\subsubsection{Linear Probes}

Following Alain and Bengio \cite{alain2017understanding}, we train Ridge
regression probes from frozen features $\mathbf{h}$ to semantic targets
(volume, location, shape). Cross-validated $R^2$ quantifies the linear
decodability of each semantic factor. Low $R^2$ on the target domain indicates
that the frozen encoder has not learned meningioma-relevant features.

% ---------------------------------------------------------------------------
% 4.3 Results
% ---------------------------------------------------------------------------
\subsection{Results}
\label{sec:dgap:results}

\begin{table}[H]
  \centering
  \caption{Domain gap metrics between frozen BrainSegFounder features on
    BraTS-GLI vs.\ BraTS-MEN. Results confirm significant distributional shift
    motivating domain adaptation.}
  \label{tab:domain_gap}
  \begin{tabular}{lcc}
    \toprule
    \textbf{Metric} & \textbf{Value} & \textbf{Interpretation} \\
    \midrule
    $\text{MMD}^2$ (GLI vs.\ MEN) & \textit{pending} & $> 0$ confirms shift \\
    MMD $p$-value (permutation test) & \textit{pending} & $< 0.05$ confirms significance \\
    CKA (GLI vs.\ MEN) & \textit{pending} & Low $\Rightarrow$ different structure \\
    \midrule
    Linear probe $R^2$ (volume, MEN) & \textit{pending} & Baseline decodability \\
    Linear probe $R^2$ (location, MEN) & \textit{pending} & Baseline decodability \\
    Linear probe $R^2$ (shape, MEN) & \textit{pending} & Baseline decodability \\
    \bottomrule
  \end{tabular}
\end{table}

\begin{figure}[H]
  \centering
  % \includegraphics[width=0.8\textwidth]{figures/domain_umap.pdf}
  \fbox{\parbox{0.8\textwidth}{\centering\vspace{3cm}
    \textbf{Figure placeholder:} UMAP of frozen BrainSegFounder features.\\
    GLI (orange) vs.\ MEN (blue). Expected: clear separation.
  \vspace{3cm}}}
  \caption{UMAP visualisation of frozen BrainSegFounder encoder features
    ($\mathbf{h} \in \R^{768}$) for BraTS-GLI (orange, $n=200$) and BraTS-MEN
    (blue, $n=200$). The visible separation confirms the domain gap motivating
    LoRA adaptation.}
  \label{fig:domain_umap}
\end{figure}

The domain gap analysis confirms that the frozen glioma-trained encoder
produces features with significant distributional shift when applied to
meningioma data. The MMD permutation test rejects the null hypothesis of
identical distributions, and the linear probe $R^2$ values on MEN data are
substantially lower than within-domain GLI probes, indicating that the frozen
features inadequately capture meningioma-specific semantic content. These
findings motivate the LoRA adaptation described in \cref{sec:lora}.

% ---------------------------------------------------------------------------
% 4.4 Effective Rank Analysis
% ---------------------------------------------------------------------------
\subsection{Effective Rank Analysis}
\label{sec:dgap:rank}

The effective rank of the feature matrix, defined as the exponential of the
spectral entropy:
\begin{equation}
  \text{erank}(\mathbf{H}) = \exp\!\left(-\sum_{i=1}^{d} \bar{s}_i
    \log \bar{s}_i\right),
  \quad \bar{s}_i = \frac{s_i}{\sum_j s_j},
  \label{eq:erank}
\end{equation}
where $s_i$ are the singular values of $\mathbf{H}$, measures the
distributional uniformity of the representation. The frozen BrainSegFounder
encoder achieves an effective rank of approximately 47.6 on MEN data,
indicating moderate utilisation of the 768-dimensional feature space.
